\begin{table}
\caption{Supplementary Table S5. Support behind each RQ3 column. The primary analysis requires 20 positives and 20 negatives per diagnosis in both strata; the final two columns repeat it at a threshold of 10, which is reported because the threshold changes how many held-out sources can contribute. Structural and rhythm defects are rare outside Ningbo, so their contrasts rest on fewer sources than the spectral indicators, and the effect sizes should be read with that in mind.}
\begin{tabular}{lrrrrrlrr}
\toprule
 & dAUPRC & lo & hi & held-out sources & source x diagnosis contrasts & sources & dAUPRC (support 10) & sources (support 10) \\
indicator &  &  &  &  &  &  &  &  \\
\midrule
HF noise & 0.052300 & 0.041700 & 0.064100 & 4 & 46 & Chapman,Georgia,Ningbo,PTBXL & 0.051200 & 4 \\
baseline wander & 0.010300 & -0.001100 & 0.022300 & 4 & 51 & Chapman,Georgia,Ningbo,PTBXL & 0.010300 & 4 \\
mains & 0.056300 & 0.045200 & 0.066900 & 4 & 46 & Chapman,Georgia,Ningbo,PTBXL & 0.048500 & 4 \\
flat/clipped lead & -0.149700 & -0.190500 & -0.108300 & 1 & 5 & Ningbo & -0.116400 & 1 \\
QRS failure & -0.098500 & -0.129800 & -0.068200 & 2 & 16 & Ningbo,PTBXL & -0.083100 & 4 \\
RR implausibility & -0.126800 & -0.150600 & -0.102900 & 2 & 20 & Ningbo,PTBXL & -0.114100 & 4 \\
\bottomrule
\end{tabular}
\end{table}
